% =====================================================================
% intro_linkage.tex
%
% Linkage between the introduction of PAPER_B_MANUSCRIPT.tex and the rest
% of the paper. Each introduction paragraph gets a one-word name, and every
% unit of the body is tagged with the introduction paragraph(s) it expands.
% Two checks follow: an introduction paragraph with no body unit tagged to
% it is a promise the paper does not keep, and a body unit with no tag is
% content the introduction never announces.
%
% The introduction is mapped as it stands at ff1d52ff, in six paragraphs.
% Paragraph 2 now ends with the stance that the plan had given a paragraph
% 3A of its own, so Stance is a job hosted at the end of paragraph 2 and
% Channels is paragraph 3. The roadmap paragraph is excluded, since every
% section links to it trivially.
%
% The paragraph jobs were drafted on 2026-10-08 from the headline, which
% the author settled first; each job expands one stretch of the headline.
% On 2026-10-09 each job was set beside what its paragraph now says, the
% introduction homes in the body mapping were rechecked, and the report
% gives each earlier item its status. Only the introduction changed
% between b3b12e42 and 02696a3d, so the body units are as before. Later
% the same day bda3c559 edited the abstract and paragraph 2, and ff1d52ff
% cut Section 3's Heckman audit paragraph and Bohren partiality paragraph,
% whose rows are removed. d03e95ad and e3a4e529 cut the point that an
% audit's answer turns on the statistic rather than on how many evaluators
% are averaged from paragraph 6, Section 5 and the conclusion.
%
% A unit is a paragraph, or a numbered result together with the prose and
% proof pointer around it. Appendix proofs are mapped to the result they
% prove. Declarations and acknowledgements are not mapped.
% =====================================================================
\documentclass[11pt,a4paper]{article}
\usepackage[margin=2.2cm]{geometry}
\usepackage{amsmath,amssymb}
\usepackage{booktabs}
\usepackage{array}
\usepackage{longtable}

\newcommand{\none}{\textbf{none}}
\newcommand{\PB}{P^{\mathrm{B}}}
% Drafts below are meant for pasting into the manuscript, so citations print their keys.
\newcommand{\citep}[1]{\texttt{[#1]}}
\sloppy

\title{Introduction linkage: what each introduction paragraph promises,\\ and where the paper keeps it}
\author{}
\date{}

\begin{document}
\maketitle

\noindent Each paragraph of the introduction is named by its job, and each job is set beside
what the paragraph says at ff1d52ff (Table~\ref{tab:names}). Every unit of the body is then
tagged with the introduction paragraphs it expands (Table~\ref{tab:map}). An introduction
paragraph with few body units behind it is thinly expanded, and a body unit tagged \none\ is
content the introduction does not announce (Section~\ref{sec:findings}). Paragraph 2 hosts two
jobs, Literature and Stance, and the roadmap is left out. The jobs are drafted from the headline
below, so that the paragraphs read in order compact into it.

\paragraph{Headline.} The introduction's paragraphs compacted into one sentence, in their order.
\begin{quote}
Sequence effects on judgment are well established but not how far they can be measured, and
Hawthorne leaves open whether evaluators adopt each impression in full or adjust their belief by
a factor; the paper does not resolve that question but measures how full adoption departs from
the sequence-free factor reading of the same impressions, a departure that exists only when the
attributes are believed related and that, for panels reading a credential and a letter in
opposite sequences, is first order in the marginals and the share of decisions changed but
second order in the believed association and the surplus lost; so at a given precision the
sequence-free reading is indistinguishable from sequence-dependent full adoption through some
questions and distinguishable through others, and the paper says exactly which smooth
statistics of the average belief fall on each side.
\end{quote}

\small
\begin{longtable}{@{}p{0.7cm}p{1.6cm}p{6.5cm}p{6.5cm}@{}}
\caption{The introduction's paragraphs, named by their job, and what each says at ff1d52ff.}\label{tab:names}\\
\toprule
\textbf{\P} & \textbf{Name} & \textbf{Job} & \textbf{At ff1d52ff}\\
\midrule
\endfirsthead
\toprule
\textbf{\P} & \textbf{Name} & \textbf{Job} & \textbf{At ff1d52ff}\\
\midrule
\endhead
\bottomrule
\endfoot
1 & Gap & Sequence effects on judgment are long established, but how far they can be measured, in the statistics of beliefs and decisions an observer reads, has had little attention.
  & States the gap, and previews the share of decisions changed at first order and the association and the surplus lost at second. The marginals are not named. The second sentence argues the gap's premise from personal histories (item 11).\\
\addlinespace
2 & Literature & Work on Jeffrey conditioning asks whether updates in opposite sequences agree, not how far their disagreement can be measured; Hawthorne's factor inputs remove the dependence but leave open whether an observer can tell which rule evaluators follow.
  & As the job. Its last literature sentence says the factor inputs leave open whether an observer can tell the rules apart (item 10, resolved).\\
\addlinespace
2, end & Stance & Hawthorne leaves open whether evaluators adopt each impression in full or adjust their belief by a factor; the paper does not resolve that question but measures how full adoption departs from the sequence-free factor reading of the same impressions, which is the benchmark.
  & Hawthorne's question, then one sentence that sets full adoption against partial adoption and measures how far full adoption, ``which Hawthorne finds unrealistic'', can depart from the sequence-free factor reading (item 17). The factor reading is not called the benchmark here.\\
\addlinespace
3 & Channels & The sequence acts through what the evaluator believes, the association channel ($c$), or through the weight the evaluator gives a cue by its place in the sequence, the position channel ($\omega$); under full adoption ($\omega=1$) only the association channel acts, so the departure exists only when the attributes are believed related.
  & As the job, with the position channel defined as a weight by place (02696a3d). Adds the three coordinates of a belief and that a difference below the precision is invisible. The sentence giving the orders was cut in 3fdcac51.\\
\addlinespace
4 & Vignette & Two hiring panels read a credential and a reference letter in opposite sequences; when the traits are believed related, each panel's belief departs from the factor reading at first order in the prevalence of the trait it reads last and at second order in the association, and across many files the share of hiring decisions changed is first order and the surplus lost second order.
  & Two panels read the file in opposite sequences and, under full adoption, differ only when the traits are believed related. The rest sets full adoption against partial adoption and points to Section~6. None of the job's orders is stated (item 16).\\
\addlinespace
5 & Precision & ``Order'' is order in $c$ and ``sequence'' is reading order; an observer resolves a departure only when it exceeds the precision $\varepsilon$, so with $c^{2}\ll\varepsilon\ll c$ the factor reading is indistinguishable from full adoption through the second-order questions and distinguishable through the first-order ones.
  & As the job. The Domotor sentence (item 3) and ``the two sequences separate'' and ``arrival order'' (item 13) remain.\\
\addlinespace
6 & Results & The vignette's orders are proved in the model (DIV, DRF, SHR at first order; DEC, LOS at second), and a smooth statistic of the average belief is second order exactly when its differential at independence is proportional to the association's (PRO), which says on which side of the precision each such question falls.
  & DIV, DEC, LOS and PRO, then the belief-audit sentence (item 14). DRF and SHR are not named (item 15).\\
\end{longtable}

\begin{longtable}{@{}p{1.7cm}p{9.6cm}p{4.0cm}@{}}
\caption{The body of the paper mapped to the introduction at ff1d52ff.}\label{tab:map}\\
\toprule
\textbf{Location} & \textbf{Content} & \textbf{Intro home}\\
\midrule
\endfirsthead
\toprule
\textbf{Location} & \textbf{Content} & \textbf{Intro home}\\
\midrule
\endhead
\bottomrule
\endfoot
\S2.1 & Belief as a joint law; the association defined; the prior as the evaluator's credence; $\alpha$, $\beta$, $c$ as beliefs, not frequencies & Channels\\
\S2.1 & Impressions are attribute-local credences, not likelihoods; the sequence $\sigma$, the share $\lambda$, the adoption weight $\omega$ & Stance, Channels, Vignette\\
\S2.1 & The engagement decision (desirability $v$, threshold $\tau$, score) & Gap, Results (both weak)\\
\S2.1 & The surplus defined at the benchmark; the loss $L(c)$ & Gap, Results\\
\addlinespace
\S2.2 & Jeffrey's rule, Diaconis--Zabell's support for it, the sentence placing all results at $\omega=1$ & Stance, Channels\\
\S2.2 & Definition of the sequence effect & Vignette\\
\S2.2 & Definition of the gap; the benchmark as the same impressions read as Bayes factors & Stance (see item 12)\\
\S2.2 & Example 1, the panel's numbers & Vignette\\
\addlinespace
\S2.3 & Assumption 1, small $c$ & Precision\\
\S2.3 & Assumption 2, soft cues & Vignette, Stance\\
\S2.3 & Assumption 3, attribute-locality & Channels\\
\S2.3 & Assumption 4, surplus criterion & Gap, Results\\
\addlinespace
\S3 & The Döring and Hawthorne debate; the two responses (a defect to repair, the input re-described) & Literature, Stance\\
\S3 & The information-theoretic ground for the premise; the degree of adoption as something to measure & Stance\\
\S3 & Observational learning (Bikhchandani--Hirshleifer--Welch, Banerjee) and the dynamics of discrimination & \none\ (item 2)\\
\S3 & Decision-theory axiomatisations, Pettigrew--Weisberg, cumulative prospect theory & \none\ (item 2)\\
\S3 & Statistical discrimination, Coate--Loury, the referral gap, stereotypes as selective recall & Results (weak; item 2)\\
\addlinespace
\S4 opening & First against second order; why these statistics (three coordinates of a belief; how many against how much) & Channels, Precision, Results\\
\S4.1 & Proposition IMM, immunity on a single cue and at $c=0$ & Channels, Vignette\\
\S4.1 & Proposition DIV, its mechanism and conditions; the individual sequence effect & Results, Channels\\
\S4.1 & The score gap & Results\\
\S4.2 & Lemma ASC; the marginal as a read-out of the gap & Vignette, Precision, Results\\
\S4.3 & The flip indicator as a step function; the individual loss & Results\\
\addlinespace
\S5 opening & The classification survives averaging; distortion and protection defined & Results, Precision\\
\S5 & Proposition ORD, the sequence effect read through the statistic, with no benchmark & Results, Vignette\\
\S5 & The answer is settled before averaging; the precision argument ($\varepsilon$) and the panel reviewer & Precision, Results\\
\S5 & Table 1 & Results\\
\S5.1 & Aggregate gap and score gap; Lemma SCR & Results\\
\S5.2 opening & Averaging leaves each statistic's individual order as it was; the statistic, not the averaging, decides & Results\\
\S5.2 & The stereotype defined; Lemma SEP for any number of attributes; Proposition DEC & Results, Channels\\
\S5.2 & The odds ratio unchanged by every separable reweighting & \none\ (item 5)\\
\S5.2 & The vignette's implication: a reviewer of the average association sees nothing & Vignette, Precision\\
\S5.2 & Proposition DRF and its consequence & Precision (weak; item 15)\\
\S5.3 opening & Share and loss read the same threshold band, counted against weighted by stake & Gap, Results\\
\S5.3 & Theorem LOS, the band figure, Proposition SHR, the audit implication & Gap, Results (SHR in Gap only; item 15)\\
\S5.4 & Proposition PRO, the protected class, the panel's audit questions & Results\\
\addlinespace
\S6 & The assumptions revisited & Precision, Stance, Channels, Results\\
\S6 & The homogeneous prior: every evaluator shares one $c$ & \none\ (item 6)\\
\S6 & Partial adoption, the four cases, Proposition LAD, the cases table & Vignette, Channels, Stance\\
\S6 & The Hawthorne paragraph & Stance\\
\addlinespace
\S7 & The audit-questions table; the central contribution; the two implications for audits & Gap, Results\\
\S7 & Locating the objection in the levels rather than in the believed link; the normative question left where Döring and Hawthorne leave it & Gap, Stance, Literature\\
\addlinespace
App.~A & Proofs of DIV, SEP, LOS and PRO & Results\\
App.~A & Proof of LAD & Channels\\
\end{longtable}
\normalsize

\section{Report}\label{sec:findings}

Items 1 to 14 come from the mapping of 2026-10-08 and keep their numbers. Their status is given
at ff1d52ff. Items 15 to 17 are new. Reading the introduction against the headline found two
errors not listed before, the argument returned to paragraph 1 under item 11 and the wording
attributed to Hawthorne under item 17. Drafts
follow \texttt{notes/writing\_discipline.md}.

\subsection*{Resolved}

\paragraph{1. Gap: ``any number of cues and attributes''.} The sentence was cut from paragraph 1
in 3fdcac51, and the vignette's clause on more cues in 6668d98c. The general case of Lemma SEP now
has no introduction statement, which the introduction need not carry.

\paragraph{7. The three coordinates of a belief.} Paragraph 3 now names them (6668d98c), and
paragraph 1 previews the two decision statistics.

\paragraph{10. Paragraph 2's ``cannot tell''.} The sentence now says the factor inputs leave open
whether an observer can tell the rules apart (6668d98c).

\subsection*{Errors still open}

\paragraph{3. Precision: the Domotor sentence.} The ``myopic structure'' sentence was cut
(02696a3d). The Domotor sentence remains and cites a source that does not support the mechanism,
so plan entry 1.7 (C.7) still applies to it. That entry's other half, the typos in the terminology
sentence, is now in the manuscript, with ``proximity to'' where the entry has ``proximity
with''.

\paragraph{9. ``Correlated''.} The abstract's ``two correlated soft cues'' is gone (bda3c559).
Three places still call the attributes or the cues correlated, where the model has an evaluator
who believes the attributes related. They are Assumption 1 (``the two correlated attributes'',
found on 2026-10-09), Assumption 2 (``bundled, correlated, and soft''), and paragraph 4 (``even if
the traits were uncorrelated'', proposal 4b of 2026-10-09, not taken). Plan entry 2.6 rewrites
Assumption 2.

\paragraph{11. Paragraph 1 argues the premise from personal histories.} The ``unlikely
situation'' sentence was cut in 6668d98c, and its replacement makes a related argument. It
reads ``if individual judgments did not depend on the cue-arrival sequence at all, then
subjective experiences should not depend on personal histories either''. Under the factor
reading the belief depends on every impression received and is free only of their sequence, so
judgments free of sequence still depend on personal histories, and the inference does not hold.
The sentence also judges the factor rule on plausibility, the ground the paper sets aside.
\emph{Remedy (Gap).} Cut the clause from ``since'' to the end of the sentence. The first
sentence's citations already carry that sequence effects are well established.

\paragraph{12. Setup 2.2 says the benchmark is not a rival rule.} Unchanged at line 199, ``The
benchmark $\PB$ is not a rival updating rule but the same two impressions re-expressed''. The
headline sets the factor reading against full adoption as the rule Hawthorne prefers.
\emph{Remedy (Setup 2.2).}
\begin{quote}
The benchmark $\PB$ is the same two impressions re-expressed, each read as a multiplicative factor
relative to the prior marginal.
\end{quote}

\paragraph{13. ``Order'' for reading sequence.} Three places remain. Paragraph 5's ``the two
sequences separate'' and ``structurally blind to arrival order'' were proposals 5d and 5f of
2026-10-09, not taken. Section 5.2's ``whether arrival order matters or not'' (line 715) should
read ``arrival sequence''.

\paragraph{14. Typos and one misstatement in paragraph 6.} ``bechmark'', the backquotes closing
``sequence'' and ``while evaluator diverges'' are fixed (02696a3d), and ``and audit's answer'' went
with its sentence (d03e95ad). ``One consequence, is'' and ``may not even be detectable'' remain.
Plan entry 1.8 (C.8) covers them, and was proposal 6b of 2026-10-09, not taken.

\paragraph{17. Paragraph 2: ``which Hawthorne finds unrealistic''.} Added in bda3c559.
Hawthorne's word is ``implausible'' (``highly implausible'', p.~97; ``it seems implausible
that\ldots'', pp.~98--99), and at p.~115 ``psychologically less plausible'', all verified in the
citation audit (H3, H6, H7). The parenthesis also sits in the sentence that sets full adoption
against partial adoption, so a reader can take partial adoption as Hawthorne's alternative. The
audit found that attribution unsupported (H34), since his alternatives are factor models.
\emph{Remedy (Stance).} ``(which Hawthorne finds implausible, pp.~98--99)''.

\subsection*{Open, incomplete rather than wrong}

\paragraph{2. Literature: only the Jeffrey and belief-adjustment strand.}
\emph{Finding.} Paragraph 2 previews the commutation debate, belief adjustment and Hawthorne's
factor models. Most of Section 3 positions the paper against economics (observational learning,
decision-theory axiomatisations and statistical discrimination), and the introduction
never says so.
\emph{Remedy (Literature).} One sentence before the stance at the end of paragraph 2:
\begin{quote}
In economics, sequence has been studied across agents, as in observational learning
\citep{BHW1992,Banerjee1992}, and discrimination through evaluators' beliefs and preferences
about groups \citep{Phelps1972,Arrow1973,Bohren2019}, while decision theory axiomatises updating
on a single piece of news \citep{Epstein2006,Ortoleva2012}, and none of these asks how one
evaluator's reading sequence registers in what an audit measures.
\end{quote}
\emph{Grounds.} Each description is Section 3's own; the last clause is the paper's positioning
claim, as Section 3 states it (a single agent, two cues, a kinematic rather than an equilibrium
object).

\paragraph{4. Section 3's economics strands.} Covered by item 2.

\paragraph{5. The odds-ratio invariance.}
\emph{Finding.} Section 5.2 shows the odds ratio is unchanged by every separable reweighting, so
the odds ratio of each evaluator's belief is the same under both sequences and the benchmark, at
any precision. It is the one question that never tells the two rules apart, and the introduction
does not mention it.
\emph{Remedy (Precision).} A clause at the end of paragraph 5's precision argument:
\begin{quote}
\ldots, and the odds ratio of each evaluator's belief does not tell them apart at any precision
(Lemma~SEP).
\end{quote}
\emph{Grounds.} PropDEC.lean \texttt{oddsRatio\_gap\_eq\_zero}; Ladder.lean
\texttt{oddsRatio\_rescale}. The claim is for each evaluator's belief and not for the average
belief, whose odds ratio is a smooth statistic covered by PRO.

\paragraph{6. The homogeneous prior.}
\emph{Finding.} Section 6 states that every evaluator shares one prior $(\alpha,\beta,c)$ and that a
population differing in $c$ is outside the results. The introduction's audit claims are claims
about populations, and they carry this condition silently.
\emph{Remedy (Results).} A qualifier in paragraph 6's PRO sentence: ``for every interior mixture of
reading sequences in a population whose members share one prior, and across an open set of
priors''. The alternative is to leave scope limits to Section 6, which the roadmap points to.

\paragraph{8. The decision model.}
\emph{Finding.} Paragraph 1 speaks of decisions changed by the sequence, but the introduction
never says how a belief becomes a decision.
\emph{Remedy (Vignette).} One sentence in paragraph 4, after the panels' differing beliefs:
\begin{quote}
The panel hires when the expected value of the candidate under its belief exceeds the value of
its outside option, so a change of belief changes the decision only for candidates near that
threshold.
\end{quote}
\emph{Grounds.} Setup 2.1 (engage if and only if $s(P)\ge\tau$); Decision.lean \texttt{flips\_iff},
\texttt{not\_flips\_of\_pointing\_away}.

\paragraph{15. The marginals' first order has no introduction statement.}
\emph{Finding.} Paragraph 3's sentence giving the marginals and the share at first order was cut
in 3fdcac51. Paragraph 1 previews the share at first order and does not name the marginals.
Paragraph 5 says only that the marginals' sequence dependence ``follows easily''. Paragraph 6
names DIV, DEC, LOS and PRO and does not name DRF or SHR. The headline's ``first order in the
marginals'' is therefore carried by the abstract and by no introduction paragraph.
\emph{Remedy (Results).} The first part of paragraph 6 could name both:
\begin{quote}
First, we establish that each evaluator's belief diverges from the benchmark at first order in
$c$ (Proposition~DIV), and so do at least one marginal of the average belief
(Proposition~DRF) and the share of decisions the sequence changes (Proposition~SHR), while the
average believed association between attributes (Proposition~DEC) and the surplus-weighted loss
of the population's decisions (Theorem~LOS) are second order.
\end{quote}
\emph{Grounds.} DRF ends ``For every $\lambda\in[0,1]$, therefore, at least one marginal drifts at
first order''; SHR gives the share as $f(0)\cdot\mathcal{O}(c)$. Records as in
\texttt{verification\_coverage.md}; SHR's exact constant is a sympy record only.

\paragraph{16. Paragraph 4's job moved.}
\emph{Finding.} The vignette now carries the adoption contrast and points to Section 6, so
Section 6's partial adoption, cases table and Proposition LAD gain a home in paragraph 4. The
orders that the job gave the vignette are announced by paragraphs 1, 5 and 6 instead, and the
body rows that expand them (Section 4.2, the reviewer in Section 5.2, the opening of Section 5.3)
are re-tagged in Table~\ref{tab:map}. No error follows. The cost is item 15, since the marginals
were the order the vignette alone had carried.

\end{document}
